% c12 —— §5.4 电荷共轭变换与行列式表示；§5.5 Harish-Chandra 定理与 $\tau$ 函数；参考文献
% （原书 p.84–89；p.84 页首两行已由 c11 收尾译完）
% 负责：trans-f

% ===== 原书 p.84 =====
\subsection{电荷共轭变换与行列式表示}

现在，对 $f=f(\partial^{-1})\in\C((\partial^{-1}))$，把 $\partial^{-1}$ 换成 $u$ 所得的 $\C((u))$ 的元素 $f(u)$，称为 $f(\partial^{-1})$ 的电荷共轭变换。（名目上很夸张，实质上几乎什么也没做。）下面用 $\{f_i(u)\}$ 写出 $\tau$ 函数。这里设 $f_i\in u^{-i}+u^{-i+1}\C[[u]]$（$i=1,2,\dots$）。

\begin{lemma}
\[
\tau_m=
\frac{\det
\begin{pmatrix}
f_m(\lambda_1) & \cdots & f_1(\lambda_1)\\
\vdots & & \vdots\\
f_m(\lambda_m) & \cdots & f_1(\lambda_m)
\end{pmatrix}}
{(\lambda_1\cdots\lambda_m)^{-1}\prod_{i<j}(\lambda_i^{-1}-\lambda_j^{-1})}
\]
如此置定时，把 $\tau_m$ 经变换 $t_i=\frac{1}{i}\sum_{j=1}^m\lambda_j^i$ 看作 $t$ 的函数，再用 $t$ 的 Schur 函数 $S_Y$ 展开，则当 $Y$ 的深度不超过 $m$ 时，其系数与 $\xi_Y$ 一致。
\end{lemma}

% ===== 原书 p.85 =====
\begin{proof}
置 $f(u)=\sum_j a_{ij}u^j$。

\transnote{原书此处作 $f(u)$，前文记作 $\{f_i(u)\}$，据上下文疑应为 $f_i(u)$，译文照原文排印。}

\[
\begin{aligned}
\tau_m&=
\frac{\det
\begin{pmatrix}
\lambda_1^m f_m(\lambda_1) & \cdots & \lambda_1^m f_1(\lambda_1)\\
\vdots & & \vdots\\
\lambda_m^m f_m(\lambda_m) & \cdots & \lambda_m^m f_1(\lambda_m)
\end{pmatrix}}
{\prod_{i<j}(\lambda_i-\lambda_j)}\\
&=\det
\begin{pmatrix}
1 & \lambda_1 & \cdots & \lambda_1^m & \cdots\\
\vdots & & & & \vdots\\
1 & \lambda_m & \cdots & \lambda_m^m & \cdots
\end{pmatrix}
\begin{pmatrix}
a_{m,-m} & \cdots & a_{1,-m}\\
\vdots & & \vdots\\
a_{m,0} & \cdots & a_{1,0}
\end{pmatrix}\\
&=\sum_{-m\le i_m<\cdots<i_{-1}}
\frac{\det
\begin{pmatrix}
\lambda_1^{i_{-m}+m} & \cdots & \lambda_1^{i_{-1}+m}\\
\vdots & & \vdots\\
\lambda_m^{i_{-m}+m} & \cdots & \lambda_m^{i_{-1}+m}
\end{pmatrix}}
{\prod_{i<j}(\lambda_i-\lambda_j)}\\
&\qquad\times\det
\begin{pmatrix}
a_{m,i_{-m}} & \cdots & a_{1,i_{-m}}\\
\vdots & & \vdots\\
a_{m,i_{-1}} & \cdots & a_{1,i_{-1}}
\end{pmatrix}
\end{aligned}
\]
成为如此。这里用 Schur 函数的定义
\[
S_Y(t)=
\frac{\det
\begin{pmatrix}
\lambda_1^{i_{-m}+m} & \cdots & \lambda_1^{i_{-1}+m}\\
\vdots & & \vdots\\
\lambda_m^{i_{-m}+m} & \cdots & \lambda_m^{i_{-1}+m}
\end{pmatrix}}
{\prod_{i<j}(\lambda_i-\lambda_j)}
\]
即得引理 5.1。
\end{proof}

\begin{corollary}
引理 5.1 中 $\tau_m$ 的极限 $\lim_{m\to\infty}\tau_m$，作为无限变量对称多项式 $\varprojlim\C[s_i]$（$s_i$ 为基本对称多项式）的元素是良定义的。
\end{corollary}

\subsection{Harish-Chandra 定理与 $\tau$ 函数}

这里用 Harish-Chandra 定理改写矩阵积分。

\begin{theorem}[Harish-Chandra {[HC]}]
设 $Y=\operatorname{diag}(y_i)\in\R^n$ 为对角矩阵，且 $y_i\ne y_j$（$i\ne j$）。设 $\Phi:H_n\to\R$ 为类函数，则存在只依赖于 $n$ 的常数 $c_n$，使
\[
\int_{H_n}\exp(\tr(-\sqrt{-1}YX)\Phi(X)dX
=c_nV(Y)^{-1}\int_{\R^n}\exp(\tr(-\sqrt{-1}YD)\Phi(D)V(D)dD
\]
\end{theorem}

\transnote{上式两处 $\exp($ 在原书中均未闭合，译文照原文排印；通行的写法应补出右括号，作 $\exp(\tr(\cdots))$。}

% ===== 原书 p.86 =====
成为如此。这里对 $Y=\operatorname{diag}(y_i)$，有
\[
V(y)=\prod_{i<j}(y_i-y_j)
\]
用这一定理改写矩阵积分。把 $Y$ 延拓到 $\C$ 上成为全纯函数。又省略积分区域，作形式上的变形。
\[
\Bigl\langle\exp\Bigl(\frac{t\tr(X^3)}{3}\Bigr)\Bigr\rangle
=\int\exp\Bigl(\tr\Bigl(-\frac{1}{2}\Lambda X^2+\frac{tX^3}{3}\Bigr)\Bigr)dX
\]
这里令 $X-\frac{\Lambda}{2t}=Z$ 作积分变量替换，则上式成为
\[
\begin{aligned}
&\Bigl(\int\exp\Bigl(\tr\Bigl(\frac{t}{3}Z^3-\frac{\Lambda^2Z}{4t}\Bigr)\Bigr)dZ\Bigr)\cdot\exp\Bigl(\tr\Bigl(-\frac{\Lambda^3}{12t^2}\Bigr)\Bigr)\\
&=\int\exp\Bigl(\tr\Bigl(\frac{t}{3}Z^3\Bigr)\Bigr)\exp\Bigl(-\tr\frac{\Lambda^2Z}{4t}\Bigr)dZ\cdot\exp\Bigl(\tr\Bigl(-\frac{\Lambda^3}{12t^2}\Bigr)\Bigr)
\end{aligned}
\]
由于 $\exp(\tr(\frac{t}{3}Z^3))$ 是类函数，故用 Harish-Chandra 定理，上式成为
\[
\begin{aligned}
&c_n\frac{1}{\bigl(\frac{-\sqrt{-1}}{4t}\bigr)^{n(n-1)/2}\prod_{i<j}(\lambda_i^2-\lambda_j^2)}\\
&\qquad\times\int V(D)\exp\Bigl(\tr\Bigl(\frac{tD^3}{3}-\frac{\Lambda^2D}{4t}\Bigr)\Bigr)dD\cdot\exp\Bigl(\tr\Bigl(-\frac{\Lambda^3}{12t^2}\Bigr)\Bigr)\\
&=c_n\frac{\exp\bigl(\tr\bigl(-\frac{\Lambda^3}{12t^2}\bigr)\bigr)}{\prod_{i<j}(\lambda_i^2-\lambda_j^2)}
\det\Bigl(\int x^{j-1}\exp\Bigl(\frac{tx^3}{3}-\frac{\lambda_i^2x}{4t}\Bigr)dx\Bigr)_{ij}\\
&=c_n\frac{1}{\prod_{i<j}(\lambda_i^2-\lambda_j^2)}
\det\Bigl(\int\Bigl(x+\frac{\lambda_i}{2t}\Bigr)^{j-1}
\exp\Bigl(-\frac{\lambda_i}{2}x^2+\frac{tx^3}{3}\Bigr)dx\Bigr)_{ij}
\end{aligned}
\]
成为如此。现在令
\[
f_j(x)=\int\Bigl(x+\frac{\lambda}{2t}\Bigr)^{j-1}
\exp\Bigl(-\frac{\lambda}{2}x^2+\frac{t}{3}x^3\Bigr)dx
\]
则上式成为
\[
c'_n\frac{1}{\prod_{i<j}(\lambda_i^2-\lambda_j^2)}\det(f_j(\lambda_j))_{ij}
\]

\transnote{原书此处作 $\det(f_j(\lambda_j))_{ij}$，按上文应为 $\det(f_j(\lambda_i))_{ij}$（即 $\lambda_i^2$ 与 $x+\frac{\lambda_i}{2t}$ 中的 $\lambda_i$），译文照原文排印。}

成为如此。下面看 $f_j(\lambda)$ 关于 $\lambda$ 具有怎样的展开。先展开 $\exp(\frac{t}{3}x^3)$ 的部分来考虑，则只须考虑下面的积分（$1\le k\le j-1$，$m\ge0$）。
\[
\int x^{k+3m}\lambda^{j-k-1}\exp\Bigl(-\frac{\lambda x^2}{2}\Bigr)dx
=c_{j,k,m}\lambda^{-\frac{1}{2}(-2j+3+3m+3k)}
\]

% ===== 原书 p.87 =====
其中 $c_{j,k,m}$ 可看作 $k+3m$ 跑遍 $0$ 以上的偶数。由于
\[
f_j=\lambda^{j-\frac{3}{2}}(a_{j0}+a_{j1}\lambda^{-1}+a_{j2}\lambda^{-2}+\cdots)
\]
具有这种形状，故置
\[
h_j(\lambda^{-1})=\lambda^j(a_{j0}+a_{j1}\lambda^{-1}+a_{j2}\lambda^{-2}+\cdots)
\]
则 $f_j(\lambda)=h_j(\lambda^{-1})\lambda^{-\frac{3}{2}}$。因此
\[
\int\exp\Bigl(-\frac{\Lambda X^2}{2}+\frac{tX^3}{3}\Bigr)dX
=c_n\frac{\prod_{i=1}^n(\lambda_i)^{-\frac{3}{2}}}{\prod_{i<j}(\lambda_i^2-\lambda_j^2)}
\det(h_j(\lambda_i^{-1}))
\]
又
\[
\langle1\rangle=\exp\Bigl(-\frac{\Lambda X^2}{2}\Bigr)dX
=\frac{d_n}{\prod_{i<j}(\lambda_i+\lambda_j)\cdot\prod_{i=1}^n\lambda_i^{\frac{1}{2}}}
\]

\transnote{原书此处作 $\langle1\rangle=\exp\bigl(-\frac{\Lambda X^2}{2}\bigr)dX$，缺积分号（上下文中的 $\langle\exp(\cdot)\rangle$ 均带 $\int$），译文照原文排印。}

（$d_n$ 是只依赖于 $n$ 的常数。）故
\[
\begin{aligned}
\frac{\langle\exp(\frac{t\tr(X^3)}{3})\rangle}{\langle1\rangle}
&=c_n''\frac{\prod_{i<j}(\lambda_i+\lambda_j)\prod(\lambda_i)^{\frac{1}{2}}}
{\prod_{i<j}(\lambda_i^2-\lambda_j^2)\prod(\lambda_i)^{\frac{3}{2}}}\det(h_j(\lambda_i^{-1}))\\
&=c_n''\frac{1}{\prod_{i<j}(\lambda_i-\lambda_j)\prod_{i=1}^n\lambda_i}\det(h_j(\lambda_i^{-1}))
\end{aligned}
\]
成为如此。因此，作变量替换 $\lambda_i^{-1}=\xi_i$、$t_j=\sum_{i=1}^n\xi_i^j/j$，把
\[
\frac{\langle\exp(\frac{t\tr(X^3)}{3})\rangle}{\langle1\rangle}
\]
看作 $(t_1,t_2,\dots)$ 的函数，它便是关于 $(t_1,t_2,\dots)$ 的 $\tau$ 函数。特别地，令 $t=\frac{1}{2}$，则
\[
\frac{\langle\exp(\frac{\tr X^3}{6})\rangle}{\langle1\rangle}
\]
是关于 $(t_1,t_2,\dots)$ 的 $\tau$ 函数。由此得到下述定理。

\begin{theorem}
\[
\exp\Bigl(\Bigl\langle\sum_i T_i\tau_i\Bigr\rangle\Bigr)
=\frac{\langle\exp(\frac{\tr X^3}{6})\rangle}{\langle1\rangle}
\]
经变量替换 $(2d+1)!!T_d=t_{2d+1}$，是 KP 方程的 $\tau$ 函数。
\end{theorem}

\begin{remark}
由此可知 $\frac{\langle\exp(\frac{\tr X^3}{6})\rangle}{\langle1\rangle}$ 只依赖于 $t_{2d+1}$。只依赖于 $t_{2d+1}$ 的 $\tau$ 函数，称为 KdV 方程的 $\tau$ 函数。
% ===== 原书 p.88 =====
一般地，设 $N$ 为自然数，不依赖于变量 $t_{Nk+1}$ 的 $\tau$ 函数，称为 $N$ reduce 型 $\tau$ 函数。作为上述矩阵积分的推广，考虑下面的矩阵积分。（[KS]）
\[
\int_{H_n}\exp\Bigl(\tr\Bigl(-\sqrt{-1}\Bigl(\frac{X^{d+1}}{d+1}-\Lambda X\Bigr)\Bigr)\Bigr)dX
\]
已知其渐近展开是 $N$ reduced 型 $\tau$ 函数，但这一 $\tau$ 函数在模空间意义上的含义至今不明。
\end{remark}

\transnote{原书对同一概念前后分别用 reduce 与 reduced 两种写法，译文分别照原文字面译出，所指相同。}

% ===== 原书 p.89 =====
\addcontentsline{toc}{section}{参考文献}
\section*{参考文献}
\begin{itemize}[leftmargin=2em,itemsep=3pt,topsep=4pt,label={}]
\item {}[HC]:Harish-Chandra, Differentail operators on a semi-simple Lie algebra, Amer.J. Math. 79, pp. 87-120, (1957).
\item {}[KS]:Kac,V., Schwartz, A., Geometric interpretation of the partition function of 2d gravity, Phisical Lett. B257, 329, (1991).
\item {}[L]:Looijenga, E., Intersection theory on Deligne-Mumford compactifications., Asterisque 216, Seminaire Bourbaki, pp. 187-212.
\item {}[Sa]:佐藤幹夫(野海正俊記) ソリトン方程式と普遍グラスマン多様体, 上智大学数学講究録 18.
\item {}[St]:Streble, K., Quadratic differntials, Berlin, Heiderberg, New York, Springer 1984.
\item {}[W]:Witten, E., Two dimensional gravity and intersection theory on moduli space. Surveys in Diff. Geom. 1, pp.243-310.
\end{itemize}

\transnote{原书参考文献照原文抄录，其中若干拼写照原样保留：Differentail（应为 Differential）、Phisical（应为 Physical）、differntials（应为 differentials）、Heiderberg（应为 Heidelberg）、Streble（应为 Strebel）。}
